% generated by GAPDoc2LaTeX from XML source (Frank Luebeck)
\documentclass[a4paper,11pt]{report}

\usepackage[top=37mm,bottom=37mm,left=27mm,right=27mm]{geometry}
\sloppy
\pagestyle{myheadings}
\usepackage{amssymb}
\usepackage[utf8]{inputenc}
\usepackage{makeidx}
\makeindex
\usepackage{color}
\definecolor{FireBrick}{rgb}{0.5812,0.0074,0.0083}
\definecolor{RoyalBlue}{rgb}{0.0236,0.0894,0.6179}
\definecolor{RoyalGreen}{rgb}{0.0236,0.6179,0.0894}
\definecolor{RoyalRed}{rgb}{0.6179,0.0236,0.0894}
\definecolor{LightBlue}{rgb}{0.8544,0.9511,1.0000}
\definecolor{Black}{rgb}{0.0,0.0,0.0}

\definecolor{linkColor}{rgb}{0.0,0.0,0.554}
\definecolor{citeColor}{rgb}{0.0,0.0,0.554}
\definecolor{fileColor}{rgb}{0.0,0.0,0.554}
\definecolor{urlColor}{rgb}{0.0,0.0,0.554}
\definecolor{promptColor}{rgb}{0.0,0.0,0.589}
\definecolor{brkpromptColor}{rgb}{0.589,0.0,0.0}
\definecolor{gapinputColor}{rgb}{0.589,0.0,0.0}
\definecolor{gapoutputColor}{rgb}{0.0,0.0,0.0}

%%  for a long time these were red and blue by default,
%%  now black, but keep variables to overwrite
\definecolor{FuncColor}{rgb}{0.0,0.0,0.0}
%% strange name because of pdflatex bug:
\definecolor{Chapter }{rgb}{0.0,0.0,0.0}
\definecolor{DarkOlive}{rgb}{0.1047,0.2412,0.0064}


\usepackage{fancyvrb}

\usepackage{mathptmx,helvet}
\usepackage[T1]{fontenc}
\usepackage{textcomp}


\usepackage[
            pdftex=true,
            bookmarks=true,        
            a4paper=true,
            pdftitle={Written with GAPDoc},
            pdfcreator={LaTeX with hyperref package / GAPDoc},
            colorlinks=true,
            backref=page,
            breaklinks=true,
            linkcolor=linkColor,
            citecolor=citeColor,
            filecolor=fileColor,
            urlcolor=urlColor,
            pdfpagemode={UseNone}, 
           ]{hyperref}

\newcommand{\maintitlesize}{\fontsize{50}{55}\selectfont}

% write page numbers to a .pnr log file for online help
\newwrite\pagenrlog
\immediate\openout\pagenrlog =\jobname.pnr
\immediate\write\pagenrlog{PAGENRS := [}
\newcommand{\logpage}[1]{\protect\write\pagenrlog{#1, \thepage,}}
%% were never documented, give conflicts with some additional packages

\newcommand{\GAP}{\textsf{GAP}}

%% nicer description environments, allows long labels
\usepackage{enumitem}
\setdescription{style=nextline}

%% depth of toc
\setcounter{tocdepth}{1}





%% command for ColorPrompt style examples
\newcommand{\gapprompt}[1]{\color{promptColor}{\bfseries #1}}
\newcommand{\gapbrkprompt}[1]{\color{brkpromptColor}{\bfseries #1}}
\newcommand{\gapinput}[1]{\color{gapinputColor}{#1}}


\begin{document}

\logpage{[ 0, 0, 0 ]}
\begin{titlepage}
\mbox{}\vfill

\begin{center}{\maintitlesize \textbf{ LocalNR \\
\mbox{}}}\\
\vfill

\hypersetup{pdftitle={ LocalNR }}
\markright{\scriptsize \mbox{}\hfill  LocalNR  \hfill\mbox{}}
{\Huge \textbf{ Package of local nearrings \\
\mbox{}}}\\
\vfill

{\Huge  2.1.0 \mbox{}}\\[1cm]
{ 22 May 2026 \mbox{}}\\[1cm]
\mbox{}\\[2cm]
{\Large \textbf{\strut  Iryna Raievska     \strut\mbox{}}}\\
{\Large \textbf{\strut  Maryna Raievska     \strut\mbox{}}}\\
{\Large \textbf{\strut  Yaroslav Sysak     \strut\mbox{}}}\\
\hypersetup{pdfauthor={ Iryna Raievska    ;  Maryna Raievska    ;  Yaroslav Sysak    }}
\end{center}\vfill

\mbox{}\\
{\mbox{}\\
\small \noindent \textbf{ Iryna Raievska    }  Email: \href{mailto://raeirina@imath.kiev.ua} {\texttt{raeirina@imath.kiev.ua}}\\
  Homepage: \href{https://www.imath.kiev.ua/~algebra/Raievska_I/} {\texttt{https://www.imath.kiev.ua/\texttt{\symbol{126}}algebra/Raievska{\textunderscore}I/}}\\
  Address: \begin{minipage}[t]{8cm}\noindent
 Institute of Mathematics\\
 National Academy of Sciences of Ukraine\\
 01024 Ukraine, Kyiv, 3, Tereshchenkivska st.\\
 \end{minipage}
}\\
{\mbox{}\\
\small \noindent \textbf{ Maryna Raievska    }  Email: \href{mailto://raemarina@imath.kiev.ua} {\texttt{raemarina@imath.kiev.ua}}\\
  Homepage: \href{https://www.imath.kiev.ua/~algebra/Raievska_M/} {\texttt{https://www.imath.kiev.ua/\texttt{\symbol{126}}algebra/Raievska{\textunderscore}M/}}\\
  Address: \begin{minipage}[t]{8cm}\noindent
 Institute of Mathematics\\
 National Academy of Sciences of Ukraine\\
 01024 Ukraine, Kyiv, 3, Tereshchenkivska st.\\
 \end{minipage}
}\\
{\mbox{}\\
\small \noindent \textbf{ Yaroslav Sysak    }  Email: \href{mailto://sysak@imath.kiev.ua} {\texttt{sysak@imath.kiev.ua}}\\
  Homepage: \href{https://www.imath.kiev.ua/~algebra/Sysak/} {\texttt{https://www.imath.kiev.ua/\texttt{\symbol{126}}algebra/Sysak/}}\\
  Address: \begin{minipage}[t]{8cm}\noindent
 Institute of Mathematics\\
 National Academy of Sciences of Ukraine\\
 01024 Ukraine, Kyiv, 3, Tereshchenkivska st.\\
 \end{minipage}
}\\
\end{titlepage}

\newpage\setcounter{page}{2}
\newpage

\def\contentsname{Contents\logpage{[ 0, 0, 1 ]}}

\tableofcontents
\newpage

     
\chapter{\textcolor{Chapter }{Local nearrings}}\label{Chapter_Local_nearrings}
\logpage{[ 1, 0, 0 ]}
\hyperdef{L}{X7FA640B47E64CE5F}{}
{
  

 A set $R$ with two binary operations $+$ and $\cdot$ is called a \emph{(left) nearring} if the following statements hold: 

 
\begin{enumerate}
\item  $(R,+)=R^+$ is a (not necessarily abelian) group with neutral element $0$; 
\item  $(R,\cdot)$ is a semigroup; 
\item  $x(y+z)=xy+xz$ for all $x$, $y$, $z\in R$. 
\end{enumerate}
 

 If $R$ is a nearring, then the group $R^+$ is called the \emph{additive group} of $R$. If in addition $0\cdot x=0$, then the nearring $R$ is called \emph{zero\texttt{\symbol{45}}symmetric}, and if the semigroup $(R,\cdot)$ is a monoid, i.e. it has an identity element $i$, then $R$ is a \emph{nearring with identity} $i$. In the latter case the group $R^*$ of all invertible elements of the monoid $(R,\cdot)$ is called the \emph{multiplicative group} of $R$. 

 The concepts of a subnearring and a nearring homomorphism are defined by the
same way as for rings. In particular, if $\lambda$ is a nearring homomorphism of $(R,+, \cdot)$, then its kernel $Ker \lambda$ is a subnearring of $(R,+, \cdot)$ whose additive subgroup is normal in $R^+$. 

 A subnearring $I$ of $(R,+, \cdot)$ is an ideal of $(R,+, \cdot)$ if $I = Ker \lambda$ for some $\lambda$. 

 It can simply be verified that $I$ is an ideal of $R$ if and only if its additive group $I^+$ is a normal subgroup of $R^+$ and for any elements $r$, $s\in R$ and $a\in I$ the inclusions $ra\in I$ and $(r + a)s {\ensuremath{-}} rs\in I$ hold. Main results accumulated for local nearrings can be found in the surveys \cite{MR2799412} and \cite{MR4943879}. 

 A nearring $R$ with identity is said to be \emph{local} if the set $L=R\setminus R^*$ of all non\texttt{\symbol{45}}invertible elements of $R$ is a subgroup of $R^+$. 

 It is clear that if $L$ is an ideal of $R$, then the factor nearring $R/L$ is a \emph{nearfield}. For example, every local ring $R$ is a zero\texttt{\symbol{45}}symmetric local nearring whose subgroup $L$ coincides with the Jacobson radical of $R$. Reference: \cite{MR230773}. 
\section{\textcolor{Chapter }{The local nearrings library}}\label{Chapter_Local_nearrings_Section_The_local_nearrings_library}
\logpage{[ 1, 1, 0 ]}
\hyperdef{L}{X83C317B77CB6A7D2}{}
{
  

\subsection{\textcolor{Chapter }{AdditiveGroupsOfLibraryOfLNRsOfOrder}}
\logpage{[ 1, 1, 1 ]}\nobreak
\hyperdef{L}{X816BBAC67F5CE775}{}
{\noindent\textcolor{FuncColor}{$\triangleright$\enspace\texttt{AdditiveGroupsOfLibraryOfLNRsOfOrder({\mdseries\slshape n})\index{AdditiveGroupsOfLibraryOfLNRsOfOrder@\texttt{Additive}\-\texttt{Groups}\-\texttt{Of}\-\texttt{Library}\-\texttt{Of}\-\texttt{L}\-\texttt{N}\-\texttt{Rs}\-\texttt{Of}\-\texttt{Order}}
\label{AdditiveGroupsOfLibraryOfLNRsOfOrder}
}\hfill{\scriptsize (function)}}\\
\textbf{\indent Returns:\ }
a list 



 The argument is $n$. The output is a list of additive groups of local nearrings in the library of
this package of order $n$. }

 
\begin{Verbatim}[commandchars=!@|,fontsize=\small,frame=single,label=Example]
  !gapprompt@gap>| !gapinput@List(AdditiveGroupsOfLibraryOfLNRsOfOrder(81),IdGroup);|
  [ [ 81, 1 ], [ 81, 2 ], [ 81, 3 ], [ 81, 5 ], [ 81, 6 ], [ 81, 11 ], 
    [ 81, 12 ], [ 81, 13 ], [ 81, 15 ] ]
\end{Verbatim}
 

\subsection{\textcolor{Chapter }{LibraryOfLNRsOnGroup}}
\logpage{[ 1, 1, 2 ]}\nobreak
\hyperdef{L}{X82F6B0D18723BCC7}{}
{\noindent\textcolor{FuncColor}{$\triangleright$\enspace\texttt{LibraryOfLNRsOnGroup({\mdseries\slshape G})\index{LibraryOfLNRsOnGroup@\texttt{LibraryOfLNRsOnGroup}}
\label{LibraryOfLNRsOnGroup}
}\hfill{\scriptsize (function)}}\\
\textbf{\indent Returns:\ }
a list 



 The argument is a group $G$. The output is a list of catalogue entries for the local nearrings in the
library of this package whose additive group is isomorphic to $G$. }

 The local nearrings are sorted by their multiplicative groups. 
\begin{Verbatim}[commandchars=!@|,fontsize=\small,frame=single,label=Example]
  !gapprompt@gap>| !gapinput@G:=SmallGroup(81,2);|
  <pc group of size 81 with 4 generators>
  !gapprompt@gap>| !gapinput@LibraryOfLNRsOnGroup(G);|
  [ "AllLocalNearRings(81,2,54,3)", "AllLocalNearRings(81,2,54,6)", 
    "AllLocalNearRings(81,2,54,9)", "AllLocalNearRings(81,2,54,10)", 
    "AllLocalNearRings(81,2,54,11)", "AllLocalNearRings(81,2,54,15)", 
    "AllLocalNearRings(81,2,72,14)", "AllLocalNearRings(81,2,72,19)", 
    "AllLocalNearRings(81,2,72,24)", "AllLocalNearRings(81,2,72,26)" ]
\end{Verbatim}
 

\subsection{\textcolor{Chapter }{LocalNearRing}}
\logpage{[ 1, 1, 3 ]}\nobreak
\hyperdef{L}{X7F95FD0A7F953C41}{}
{\noindent\textcolor{FuncColor}{$\triangleright$\enspace\texttt{LocalNearRing({\mdseries\slshape k, l, m, n, w})\index{LocalNearRing@\texttt{LocalNearRing}}
\label{LocalNearRing}
}\hfill{\scriptsize (operation)}}\\
\textbf{\indent Returns:\ }
a nearring 



 The arguments are $k$, $l$, $m$, $n$, $w$. The output is the $w$\texttt{\symbol{45}}th local nearring from the library of this package whose
additive group has \texttt{IdGroup} value \texttt{[k,l]} and whose multiplicative group has \texttt{IdGroup} value \texttt{[m,n]}. No validation of the arguments is performed. }

 
\begin{Verbatim}[commandchars=!@|,fontsize=\small,frame=single,label=Example]
  !gapprompt@gap>| !gapinput@L:=LocalNearRing(81,12,54,8,3);|
  ExplicitMultiplicationNearRing ( <pc group of size 81 with 
  4 generators> , multiplication )
\end{Verbatim}
 

\subsection{\textcolor{Chapter }{AllLocalNearRings}}
\logpage{[ 1, 1, 4 ]}\nobreak
\hyperdef{L}{X83C221F187A90BBE}{}
{\noindent\textcolor{FuncColor}{$\triangleright$\enspace\texttt{AllLocalNearRings({\mdseries\slshape k, l, m, n})\index{AllLocalNearRings@\texttt{AllLocalNearRings}}
\label{AllLocalNearRings}
}\hfill{\scriptsize (operation)}}\\
\textbf{\indent Returns:\ }
a list 



 The arguments are $k$, $l$, $m$, $n$. The output is the list of all local nearrings from the library of this
package whose additive group has \texttt{IdGroup} value \texttt{[k,l]} and whose multiplicative group has \texttt{IdGroup} value \texttt{[m,n]}. No validation of the arguments is performed. }

 
\begin{Verbatim}[commandchars=!@|,fontsize=\small,frame=single,label=Example]
  !gapprompt@gap>| !gapinput@L:=AllLocalNearRings(81,12,54,8);;|
  !gapprompt@gap>| !gapinput@Size(L);|
  30
\end{Verbatim}
 

\subsection{\textcolor{Chapter }{NumberLocalNearRings}}
\logpage{[ 1, 1, 5 ]}\nobreak
\hyperdef{L}{X78DCE2A07C830905}{}
{\noindent\textcolor{FuncColor}{$\triangleright$\enspace\texttt{NumberLocalNearRings({\mdseries\slshape k, l, m, n})\index{NumberLocalNearRings@\texttt{NumberLocalNearRings}}
\label{NumberLocalNearRings}
}\hfill{\scriptsize (operation)}}\\
\textbf{\indent Returns:\ }
a number 



 The arguments are $k$, $l$, $m$, $n$. The output is the number of local nearrings in the library of this package
whose additive group has \texttt{IdGroup} value \texttt{[k,l]} and whose multiplicative group has \texttt{IdGroup} value \texttt{[m,n]}. No validation of the arguments is performed. }

 
\begin{Verbatim}[commandchars=!@|,fontsize=\small,frame=single,label=Example]
  !gapprompt@gap>| !gapinput@NumberLocalNearRings(81,15,54,8);|
  10
\end{Verbatim}
 

\subsection{\textcolor{Chapter }{IsAdditiveGroupOfLibraryOfLNRs}}
\logpage{[ 1, 1, 6 ]}\nobreak
\hyperdef{L}{X8147B516875F3499}{}
{\noindent\textcolor{FuncColor}{$\triangleright$\enspace\texttt{IsAdditiveGroupOfLibraryOfLNRs({\mdseries\slshape G})\index{IsAdditiveGroupOfLibraryOfLNRs@\texttt{IsAdditiveGroupOfLibraryOfLNRs}}
\label{IsAdditiveGroupOfLibraryOfLNRs}
}\hfill{\scriptsize (function)}}\\
\textbf{\indent Returns:\ }
a boolean 



 The argument is a group $G$. The output is \texttt{true} if the library of this package contains a local nearring whose additive group
is isomorphic to $G$, and \texttt{false} otherwise. }

 
\begin{Verbatim}[commandchars=!@|,fontsize=\small,frame=single,label=Example]
  !gapprompt@gap>| !gapinput@G:=SmallGroup(25,2);|
  <pc group of size 25 with 2 generators>
  !gapprompt@gap>| !gapinput@IsAdditiveGroupOfLibraryOfLNRs(G);|
  true
  !gapprompt@gap>| !gapinput@IsAdditiveGroupOfLibraryOfLNRs(SmallGroup(81,14));|
  false
\end{Verbatim}
 }

 }

   
\chapter{\textcolor{Chapter }{Functions}}\label{Chapter_Functions}
\logpage{[ 2, 0, 0 ]}
\hyperdef{L}{X86FA580F8055B274}{}
{
  
\section{\textcolor{Chapter }{Group functions}}\label{Chapter_Functions_Section_Group_functions}
\logpage{[ 2, 1, 0 ]}
\hyperdef{L}{X86271CEC83E77A58}{}
{
  

\subsection{\textcolor{Chapter }{IsMinimalNonAbelianGroup}}
\logpage{[ 2, 1, 1 ]}\nobreak
\hyperdef{L}{X7EDBADD77E2638BD}{}
{\noindent\textcolor{FuncColor}{$\triangleright$\enspace\texttt{IsMinimalNonAbelianGroup({\mdseries\slshape G})\index{IsMinimalNonAbelianGroup@\texttt{IsMinimalNonAbelianGroup}}
\label{IsMinimalNonAbelianGroup}
}\hfill{\scriptsize (property)}}\\
\textbf{\indent Returns:\ }
a boolean 



 The argument is a group $G$. The output is \texttt{true} if $G$ is a minimal non\texttt{\symbol{45}}abelian group, otherwise the output is \texttt{false}. }

 Recall that each finite non\texttt{\symbol{45}}abelian group whose proper
subgroups are abelian is called a \emph{Miller\texttt{\symbol{45}}Moreno group} or, in other terminology, a \emph{minimal non\texttt{\symbol{45}}abelian group}. 
\begin{Verbatim}[commandchars=!@|,fontsize=\small,frame=single,label=Example]
  !gapprompt@gap>| !gapinput@H:=SmallGroup(120,4);       |
  <pc group of size 120 with 5 generators>
  !gapprompt@gap>| !gapinput@IsMinimalNonAbelianGroup(H);|
  false
  !gapprompt@gap>| !gapinput@K:=SmallGroup(16,6);|
  <pc group of size 16 with 4 generators>
  !gapprompt@gap>| !gapinput@IsMinimalNonAbelianGroup(K);|
  true
  !gapprompt@gap>| !gapinput@IsMinimalNonAbelianGroup(SmallGroup(16,8));|
  false
\end{Verbatim}
 

\subsection{\textcolor{Chapter }{IsMetacyclicPGroup}}
\logpage{[ 2, 1, 2 ]}\nobreak
\hyperdef{L}{X83C3B54B789609D4}{}
{\noindent\textcolor{FuncColor}{$\triangleright$\enspace\texttt{IsMetacyclicPGroup({\mdseries\slshape G})\index{IsMetacyclicPGroup@\texttt{IsMetacyclicPGroup}}
\label{IsMetacyclicPGroup}
}\hfill{\scriptsize (property)}}\\
\textbf{\indent Returns:\ }
a boolean 



 The argument is a group $G$. The output is \texttt{true} if $G$ is a metacyclic $p$\texttt{\symbol{45}}group, otherwise the output is \texttt{false}. }

 
\begin{Verbatim}[commandchars=!@|,fontsize=\small,frame=single,label=Example]
  !gapprompt@gap>| !gapinput@IsMetacyclicPGroup(K);|
  true
  !gapprompt@gap>| !gapinput@IsMetacyclicPGroup(SmallGroup(81,4));|
  true
  !gapprompt@gap>| !gapinput@IsMetacyclicPGroup(SmallGroup(81,15));|
  false
\end{Verbatim}
 

\subsection{\textcolor{Chapter }{EndoOrbitsOfGroup}}
\logpage{[ 2, 1, 3 ]}\nobreak
\hyperdef{L}{X81771B097A79D6BF}{}
{\noindent\textcolor{FuncColor}{$\triangleright$\enspace\texttt{EndoOrbitsOfGroup({\mdseries\slshape G})\index{EndoOrbitsOfGroup@\texttt{EndoOrbitsOfGroup}}
\label{EndoOrbitsOfGroup}
}\hfill{\scriptsize (operation)}}\\
\textbf{\indent Returns:\ }
EndoOrbitsOfGroup 



 The argument is a group $G$. The output is a list of pairs \texttt{[x,H]}, where \texttt{x} is a representative of an endo\texttt{\symbol{45}}orbit of $G$ and \texttt{H} is the set of all images of \texttt{x} under endomorphisms of $G$. }

 
\begin{Verbatim}[commandchars=!@|,fontsize=\small,frame=single,label=Example]
  !gapprompt@gap>| !gapinput@D:=SmallGroup(81,2); |
  <pc group of size 81 with 4 generators>
  !gapprompt@gap>| !gapinput@T:=EndoOrbitsOfGroup(D);;|
  !gapprompt@gap>| !gapinput@Length(T);|
  1
  !gapprompt@gap>| !gapinput@Size(T[1][2]);|
  81
\end{Verbatim}
 

\subsection{\textcolor{Chapter }{IsEndoCyclicGroup}}
\logpage{[ 2, 1, 4 ]}\nobreak
\hyperdef{L}{X784893BB7947F14D}{}
{\noindent\textcolor{FuncColor}{$\triangleright$\enspace\texttt{IsEndoCyclicGroup({\mdseries\slshape G})\index{IsEndoCyclicGroup@\texttt{IsEndoCyclicGroup}}
\label{IsEndoCyclicGroup}
}\hfill{\scriptsize (property)}}\\
\textbf{\indent Returns:\ }
a boolean 



 The argument is a group $G$. The output is \texttt{true} if $G$ is an endocyclic group, otherwise the output is \texttt{false}. }

 Let $G$ be a group and $End G$ be the set of all its endomorphisms, which can be considered as a semigroup
with respect to composition. For each $g\in G$ we denote by $g^{End G}$ the set $\{g^\alpha| \alpha\in End G\}$ of all images of the element $g$ with respect to endomorphisms of $End G$. A group $G$ is called \emph{endocyclic} if it contains an element $g$ with $G=g^{End G}$. 
\begin{Verbatim}[commandchars=!@|,fontsize=\small,frame=single,label=Example]
  !gapprompt@gap>| !gapinput@IsEndoCyclicGroup(D);|
  true
\end{Verbatim}
 }

 
\section{\textcolor{Chapter }{Nearring functions}}\label{Chapter_Functions_Section_Nearring_functions}
\logpage{[ 2, 2, 0 ]}
\hyperdef{L}{X86A9AE187AF165C9}{}
{
  

\subsection{\textcolor{Chapter }{UnitsOfNearRing}}
\logpage{[ 2, 2, 1 ]}\nobreak
\hyperdef{L}{X877D0B1B8729395F}{}
{\noindent\textcolor{FuncColor}{$\triangleright$\enspace\texttt{UnitsOfNearRing({\mdseries\slshape R})\index{UnitsOfNearRing@\texttt{UnitsOfNearRing}}
\label{UnitsOfNearRing}
}\hfill{\scriptsize (attribute)}}\\
\textbf{\indent Returns:\ }
a list 



 The argument is a nearring $R$. The output is a list of units if $R$ is a nearring with identity, otherwise the output is an empty list. }

 
\begin{Verbatim}[commandchars=!@|,fontsize=\small,frame=single,label=Example]
  !gapprompt@gap>| !gapinput@N:=LocalNearRing(32,5,16,3,8);|
  ExplicitMultiplicationNearRing ( <pc group of size 32 with 
  5 generators> , multiplication )
  !gapprompt@gap>| !gapinput@U:=UnitsOfNearRing(N);|
  [ (f1), (f1*f5), (f1*f4), (f1*f4*f5), (f1*f3), (f1*f3*f5), (f1*f3*f4),  
   (f1*f3*f4*f5), (f1*f2), (f1*f2*f5), (f1*f2*f4), (f1*f2*f4*f5), (f1*f2*f3), 
   (f1*f2*f3*f5), (f1*f2*f3*f4), (f1*f2*f3*f4*f5) ]
  !gapprompt@gap>| !gapinput@Un:=NearRingUnits(N);;|
  !gapprompt@gap>| !gapinput@U=Un;|
  true
  !gapprompt@gap>| !gapinput@L:=LibraryNearRing(SmallGroup(6,1),3);|
  #I  using isomorphic copy of the group
  LibraryNearRing(6/2, 3)
  !gapprompt@gap>| !gapinput@UnitsOfNearRing(L);|
  [  ]
\end{Verbatim}
 

\subsection{\textcolor{Chapter }{IsLocalNearRing}}
\logpage{[ 2, 2, 2 ]}\nobreak
\hyperdef{L}{X7E2BB4F27E6B71E6}{}
{\noindent\textcolor{FuncColor}{$\triangleright$\enspace\texttt{IsLocalNearRing({\mdseries\slshape R})\index{IsLocalNearRing@\texttt{IsLocalNearRing}}
\label{IsLocalNearRing}
}\hfill{\scriptsize (property)}}\\
\textbf{\indent Returns:\ }
a boolean 



 The argument is a nearring $R$. The output is \texttt{true} if $R$ is a local nearring, otherwise the output is \texttt{false}. }

 
\begin{Verbatim}[commandchars=!@|,fontsize=\small,frame=single,label=Example]
  !gapprompt@gap>| !gapinput@H:=SmallGroup(16,6);|
  <pc group of size 16 with 4 generators>
  !gapprompt@gap>| !gapinput@A:= AutomorphismNearRing(H);|
  AutomorphismNearRing( <pc group of size 16 with 4 generators> )
  !gapprompt@gap>| !gapinput@ Size(A);|
  64
  !gapprompt@gap>| !gapinput@IsLocalNearRing(A);|
  true
  !gapprompt@gap>| !gapinput@K:=LibraryNearRingWithOne(SmallGroup(8,2),1);  |
  #I  using isomorphic copy of the group
  LibraryNearRing(8/2, 814)
  !gapprompt@gap>| !gapinput@IsLocalNearRing(K);|
  false
\end{Verbatim}
 

\subsection{\textcolor{Chapter }{IsLocalRing}}
\logpage{[ 2, 2, 3 ]}\nobreak
\hyperdef{L}{X7B37F8E185604257}{}
{\noindent\textcolor{FuncColor}{$\triangleright$\enspace\texttt{IsLocalRing({\mdseries\slshape R})\index{IsLocalRing@\texttt{IsLocalRing}}
\label{IsLocalRing}
}\hfill{\scriptsize (property)}}\\
\textbf{\indent Returns:\ }
a boolean 



 The argument is a local nearring $R$. The output is \texttt{true} if $R$ is a local ring, otherwise the output is \texttt{false}. }

 
\begin{Verbatim}[commandchars=!@|,fontsize=\small,frame=single,label=Example]
  !gapprompt@gap>| !gapinput@L:=AllLocalNearRings(16,14,8,4);;|
  !gapprompt@gap>| !gapinput@Size(L);|
  24
  !gapprompt@gap>| !gapinput@F:=Filtered(L,x->IsLocalRing(x));;|
  !gapprompt@gap>| !gapinput@Size(F);|
  1
\end{Verbatim}
 

\subsection{\textcolor{Chapter }{NearRingNonUnits}}
\logpage{[ 2, 2, 4 ]}\nobreak
\hyperdef{L}{X87400EFF857CCFD8}{}
{\noindent\textcolor{FuncColor}{$\triangleright$\enspace\texttt{NearRingNonUnits({\mdseries\slshape R})\index{NearRingNonUnits@\texttt{NearRingNonUnits}}
\label{NearRingNonUnits}
}\hfill{\scriptsize (attribute)}}\\
\textbf{\indent Returns:\ }
a set 



 The argument is a nearring $R$. The output is the set of non\texttt{\symbol{45}}invertible elements of $R$. }

 
\begin{Verbatim}[commandchars=!@|,fontsize=\small,frame=single,label=Example]
  !gapprompt@gap>| !gapinput@T:=LocalNearRing(49,2,42,1,1); |
  ExplicitMultiplicationNearRing ( <pc group of size 49 with 
  2 generators> , multiplication )
  !gapprompt@gap>| !gapinput@Nu:=NearRingNonUnits(T);|
  [ (<identity> of ...), (f2), (f2^2), (f2^3), (f2^4), (f2^5), (f2^6) ]
  !gapprompt@gap>| !gapinput@Size(Nu);|
  7
  !gapprompt@gap>| !gapinput@R:=LibraryNearRing(SmallGroup(8,4),3);|
  #I  using isomorphic copy of the group
  LibraryNearRing(8/5, 3)
  !gapprompt@gap>| !gapinput@N:=SortedList(NearRingNonUnits(R)); |
  [ (()), ((1,2,3,4)(5,6,7,8)), ((1,3)(2,4)(5,7)(6,8)), ((1,4,3,2)(5,8,7,6)), 
    ((1,5,3,7)(2,8,4,6)), ((1,6,3,8)(2,5,4,7)), ((1,7,3,5)(2,6,4,8)), 
    ((1,8,3,6)(2,7,4,5)) ]
\end{Verbatim}
 

\subsection{\textcolor{Chapter }{SubNearRingByGenerators}}
\logpage{[ 2, 2, 5 ]}\nobreak
\hyperdef{L}{X7CE616A982D407C0}{}
{\noindent\textcolor{FuncColor}{$\triangleright$\enspace\texttt{SubNearRingByGenerators({\mdseries\slshape R, gens})\index{SubNearRingByGenerators@\texttt{SubNearRingByGenerators}}
\label{SubNearRingByGenerators}
}\hfill{\scriptsize (operation)}}\\
\textbf{\indent Returns:\ }
a subnearring 



 The arguments are a nearring $R$ and generators $gens$ of $R$. The output is the subnearring generated by $gens$. }

 
\begin{Verbatim}[commandchars=!@|,fontsize=\small,frame=single,label=Example]
  !gapprompt@gap>| !gapinput@B:=LocalNearRing(25,2,20,3,1); |
  ExplicitMultiplicationNearRing ( <pc group of size 25 with 2 generators> , multiplication )
  !gapprompt@gap>| !gapinput@D:=DistributiveElements(B);;|
  !gapprompt@gap>| !gapinput@Size(D);|
  5
  !gapprompt@gap>| !gapinput@Rs:=SubNearRingByGenerators(B,D);;|
  !gapprompt@gap>| !gapinput@Size(Rs);|
  5
  !gapprompt@gap>| !gapinput@IsDgNearRing(B);|
  false
  !gapprompt@gap>| !gapinput@IsDgNearRing(Rs);|
  true
\end{Verbatim}
 

\subsection{\textcolor{Chapter }{NonUnitsAsAdditiveSubgroup}}
\logpage{[ 2, 2, 6 ]}\nobreak
\hyperdef{L}{X8321BFC6856907B6}{}
{\noindent\textcolor{FuncColor}{$\triangleright$\enspace\texttt{NonUnitsAsAdditiveSubgroup({\mdseries\slshape R})\index{NonUnitsAsAdditiveSubgroup@\texttt{NonUnitsAsAdditiveSubgroup}}
\label{NonUnitsAsAdditiveSubgroup}
}\hfill{\scriptsize (attribute)}}\\
\textbf{\indent Returns:\ }
a subgroup 



 The argument is a local nearring $R$. The output is the additive subgroup of non\texttt{\symbol{45}}units of $R$. }

 
\begin{Verbatim}[commandchars=!@|,fontsize=\small,frame=single,label=Example]
  !gapprompt@gap>| !gapinput@T:=LocalNearRing(125,4,100,9,1); |
  ExplicitMultiplicationNearRing ( <pc group of size 125 with 
  3 generators> , multiplication )
  !gapprompt@gap>| !gapinput@L:=NonUnitsAsAdditiveSubgroup(T);|
  Group([ <identity> of ..., f2, f3, f2^2, f2*f3, f3^2, f2^3, f2^2*f3, f2*f3^2, f3^3, f2^4, 
    f2^3*f3, f2^2*f3^2, f2*f3^3, f3^4, f2^4*f3, f2^3*f3^2, f2^2*f3^3, f2*f3^4, f2^4*f3^2, 
    f2^3*f3^3, f2^2*f3^4, f2^4*f3^3, f2^3*f3^4, f2^4*f3^4 ])
  !gapprompt@gap>| !gapinput@IdGroup(L);|
  [ 25, 2 ]
\end{Verbatim}
 

\subsection{\textcolor{Chapter }{NonUnitsAsNearRingIdeal}}
\logpage{[ 2, 2, 7 ]}\nobreak
\hyperdef{L}{X8113B5737BC3D587}{}
{\noindent\textcolor{FuncColor}{$\triangleright$\enspace\texttt{NonUnitsAsNearRingIdeal({\mdseries\slshape R})\index{NonUnitsAsNearRingIdeal@\texttt{NonUnitsAsNearRingIdeal}}
\label{NonUnitsAsNearRingIdeal}
}\hfill{\scriptsize (attribute)}}\\
\textbf{\indent Returns:\ }
an ideal 



 The argument is a local nearring $R$. The output is the ideal generated by all non\texttt{\symbol{45}}invertible
elements of $R$. }

 
\begin{Verbatim}[commandchars=!@|,fontsize=\small,frame=single,label=Example]
  !gapprompt@gap>| !gapinput@I:=NonUnitsAsNearRingIdeal(T);|
  < nearring ideal >
  !gapprompt@gap>| !gapinput@Size(I);|
  25
\end{Verbatim}
 

\subsection{\textcolor{Chapter }{MultiplicativeSemigroupOfNearRing}}
\logpage{[ 2, 2, 8 ]}\nobreak
\hyperdef{L}{X80619E5279778AE6}{}
{\noindent\textcolor{FuncColor}{$\triangleright$\enspace\texttt{MultiplicativeSemigroupOfNearRing({\mdseries\slshape R})\index{MultiplicativeSemigroupOfNearRing@\texttt{MultiplicativeSemigroupOfNearRing}}
\label{MultiplicativeSemigroupOfNearRing}
}\hfill{\scriptsize (attribute)}}\\
\textbf{\indent Returns:\ }
a semigroup 



 The argument is a nearring $R$. The output is the multiplicative semigroup of $R$. }

 
\begin{Verbatim}[commandchars=!@|,fontsize=\small,frame=single,label=Example]
  !gapprompt@gap>| !gapinput@B:=LocalNearRing(16,10,8,2,7);;|
  !gapprompt@gap>| !gapinput@M:=MultiplicativeSemigroupOfNearRing(B);|
  <semigroup of size 16, with 7 generators>
  !gapprompt@gap>| !gapinput@Size(M);|
  16
\end{Verbatim}
 

\subsection{\textcolor{Chapter }{NonUnitsAsMultiplicativeSemigroup}}
\logpage{[ 2, 2, 9 ]}\nobreak
\hyperdef{L}{X83280C7283BA2D92}{}
{\noindent\textcolor{FuncColor}{$\triangleright$\enspace\texttt{NonUnitsAsMultiplicativeSemigroup({\mdseries\slshape R})\index{NonUnitsAsMultiplicativeSemigroup@\texttt{NonUnitsAsMultiplicativeSemigroup}}
\label{NonUnitsAsMultiplicativeSemigroup}
}\hfill{\scriptsize (attribute)}}\\
\textbf{\indent Returns:\ }
a semigroup 



 The argument is a nearring $R$. The output is the multiplicative semigroup of non\texttt{\symbol{45}}units
of $R$. }

 
\begin{Verbatim}[commandchars=!@|,fontsize=\small,frame=single,label=Example]
  !gapprompt@gap>| !gapinput@Nm:=NonUnitsAsMultiplicativeSemigroup(B);|
  <semigroup with 8 generators>
  !gapprompt@gap>| !gapinput@Size(Nm);|
  8
\end{Verbatim}
 

\subsection{\textcolor{Chapter }{IsOneGeneratedNearRing}}
\logpage{[ 2, 2, 10 ]}\nobreak
\hyperdef{L}{X78C3A806874CB77D}{}
{\noindent\textcolor{FuncColor}{$\triangleright$\enspace\texttt{IsOneGeneratedNearRing({\mdseries\slshape R})\index{IsOneGeneratedNearRing@\texttt{IsOneGeneratedNearRing}}
\label{IsOneGeneratedNearRing}
}\hfill{\scriptsize (property)}}\\
\textbf{\indent Returns:\ }
a boolean 



 The argument is a nearring $R$. The output is \texttt{true} if $R$ is a nearring generated by one element, otherwise the output is \texttt{false}. }

 
\begin{Verbatim}[commandchars=!@|,fontsize=\small,frame=single,label=Example]
  !gapprompt@gap>| !gapinput@D:=LocalNearRing(49,2,42,4,1);|
  ExplicitMultiplicationNearRing ( <pc group of size 49 with 2 generators> , multiplication )
  !gapprompt@gap>| !gapinput@IsOneGeneratedNearRing(D);|
  true
  !gapprompt@gap>| !gapinput@H:=LocalNearRing(16,14,8,2,3);   |
  ExplicitMultiplicationNearRing ( <pc group of size 16 with 
  4 generators> , multiplication )
  !gapprompt@gap>| !gapinput@IsOneGeneratedNearRing(H);    |
  false
\end{Verbatim}
 

\subsection{\textcolor{Chapter }{AutomorphismsAssociatedWithNearRingUnits}}
\logpage{[ 2, 2, 11 ]}\nobreak
\hyperdef{L}{X836935BD79EE8810}{}
{\noindent\textcolor{FuncColor}{$\triangleright$\enspace\texttt{AutomorphismsAssociatedWithNearRingUnits({\mdseries\slshape R, Un})\index{AutomorphismsAssociatedWithNearRingUnits@\texttt{Automorphisms}\-\texttt{Associated}\-\texttt{With}\-\texttt{Near}\-\texttt{Ring}\-\texttt{Units}}
\label{AutomorphismsAssociatedWithNearRingUnits}
}\hfill{\scriptsize (operation)}}\\
\textbf{\indent Returns:\ }
a list of automorphisms 



 The arguments are a nearring $R$ with identity and a set of units $Un$ of $R$. The output is the list of automorphisms of the additive group of $R$ associated with the units in $Un$. }

 A subgroup $A$ of the automorphism group $Aut R^+$ of the additive group of the nearring $R$ with identity isomorphic to the multiplicative group $R^*$ and satisfying the condition 
\[i^A=\{i^a\mid a\in A\}=R^*\]
 is called the subgroup of $Aut R^+$ associated with $R^*$. 
\begin{Verbatim}[commandchars=!@|,fontsize=\small,frame=single,label=Example]
  !gapprompt@gap>| !gapinput@S:=UnitsOfNearRing(D);|
  [ (f1), (f1*f2), (f1*f2^2), (f1*f2^3), (f1*f2^4), (f1*f2^5), (f1*f2^6), (f1^2), (f1^2*f2), 
    (f1^2*f2^2), (f1^2*f2^3), (f1^2*f2^4), (f1^2*f2^5), (f1^2*f2^6), (f1^3), (f1^3*f2), 
    (f1^3*f2^2), (f1^3*f2^3), (f1^3*f2^4), (f1^3*f2^5), (f1^3*f2^6), (f1^4), (f1^4*f2), 
    (f1^4*f2^2), (f1^4*f2^3), (f1^4*f2^4), (f1^4*f2^5), (f1^4*f2^6), (f1^5), (f1^5*f2), 
    (f1^5*f2^2), (f1^5*f2^3), (f1^5*f2^4), (f1^5*f2^5), (f1^5*f2^6), (f1^6), (f1^6*f2), 
    (f1^6*f2^2), (f1^6*f2^3), (f1^6*f2^4), (f1^6*f2^5), (f1^6*f2^6) ]
  !gapprompt@gap>| !gapinput@A:=AutomorphismsAssociatedWithNearRingUnits(D,S);;|
  !gapprompt@gap>| !gapinput@Size(A);|
  42
\end{Verbatim}
 

\subsection{\textcolor{Chapter }{EndomorphismsAssociatedWithNearRingElements}}
\logpage{[ 2, 2, 12 ]}\nobreak
\hyperdef{L}{X847294DD81E00F8C}{}
{\noindent\textcolor{FuncColor}{$\triangleright$\enspace\texttt{EndomorphismsAssociatedWithNearRingElements({\mdseries\slshape R, Elm})\index{EndomorphismsAssociatedWithNearRingElements@\texttt{Endomorphisms}\-\texttt{Associated}\-\texttt{With}\-\texttt{Near}\-\texttt{Ring}\-\texttt{Elements}}
\label{EndomorphismsAssociatedWithNearRingElements}
}\hfill{\scriptsize (operation)}}\\
\textbf{\indent Returns:\ }
a list of endomorphisms 



 The arguments are a nearring $R$ and a set $Elm$ of nearring elements. The output is the endomorphisms associated with nearring
elements. }

 
\begin{Verbatim}[commandchars=!@|,fontsize=\small,frame=single,label=Example]
  !gapprompt@gap>| !gapinput@Nu:=NearRingNonUnits(D);|
  [ (<identity> of ...), (f2), (f2^2), (f2^3), (f2^4), (f2^5), (f2^6) ]
  !gapprompt@gap>| !gapinput@En:=EndomorphismsAssociatedWithNearRingElements(D,Nu);;|
  !gapprompt@gap>| !gapinput@Size(En);|
  7
\end{Verbatim}
 

\subsection{\textcolor{Chapter }{SemidirectProductAssociatedWithNearRing}}
\logpage{[ 2, 2, 13 ]}\nobreak
\hyperdef{L}{X7F82AB357B378109}{}
{\noindent\textcolor{FuncColor}{$\triangleright$\enspace\texttt{SemidirectProductAssociatedWithNearRing({\mdseries\slshape R})\index{SemidirectProductAssociatedWithNearRing@\texttt{Semidirect}\-\texttt{Product}\-\texttt{Associated}\-\texttt{With}\-\texttt{Near}\-\texttt{Ring}}
\label{SemidirectProductAssociatedWithNearRing}
}\hfill{\scriptsize (operation)}}\\
\textbf{\indent Returns:\ }
a semidirect product 



 The argument is a nearring $R$ with identity. The output is the semidirect product associated with the
nearring $R$. }

 
\begin{Verbatim}[commandchars=!@|,fontsize=\small,frame=single,label=Example]
  !gapprompt@gap>| !gapinput@T:=LocalNearRing(25,2,20,2,1);             |
  ExplicitMultiplicationNearRing ( <pc group of size 25 with 2 generators> , multiplication )
  !gapprompt@gap>| !gapinput@SemidirectProductAssociatedWithNearRing(T);|
  <pc group with 5 generators>
  !gapprompt@gap>| !gapinput@Size(last);|
  500
\end{Verbatim}
 

\subsection{\textcolor{Chapter }{IsCircleSubgroupOfNearRing}}
\logpage{[ 2, 2, 14 ]}\nobreak
\hyperdef{L}{X786B756C79BDA836}{}
{\noindent\textcolor{FuncColor}{$\triangleright$\enspace\texttt{IsCircleSubgroupOfNearRing({\mdseries\slshape R, H})\index{IsCircleSubgroupOfNearRing@\texttt{IsCircleSubgroupOfNearRing}}
\label{IsCircleSubgroupOfNearRing}
}\hfill{\scriptsize (operation)}}\\
\textbf{\indent Returns:\ }
a boolean 



 The arguments are a nearring $R$ with identity and a subgroup $H$ of the additive group of $R$. The output is \texttt{true} if $H$ is a circle subgroup of the nearring $R$, otherwise the output is \texttt{false}. }

 
\begin{Verbatim}[commandchars=!@|,fontsize=\small,frame=single,label=Example]
  !gapprompt@gap>| !gapinput@Sg:=Subgroups(GroupReduct(T));;|
  !gapprompt@gap>| !gapinput@Size(Sg);|
  8
  !gapprompt@gap>| !gapinput@F:=Filtered(Sg,x->IsCircleSubgroupOfNearRing(T,x));|
  [ Group([  ]), Group([ f2 ]) ]
\end{Verbatim}
 

\subsection{\textcolor{Chapter }{FactorizedGroupAssociatedWithCircleSubgroupOfNearRing}}
\logpage{[ 2, 2, 15 ]}\nobreak
\hyperdef{L}{X79C31AE57BB7CAF4}{}
{\noindent\textcolor{FuncColor}{$\triangleright$\enspace\texttt{FactorizedGroupAssociatedWithCircleSubgroupOfNearRing({\mdseries\slshape R, H})\index{FactorizedGroupAssociatedWithCircleSubgroupOfNearRing@\texttt{Factorized}\-\texttt{Group}\-\texttt{Associated}\-\texttt{With}\-\texttt{Circle}\-\texttt{Subgroup}\-\texttt{Of}\-\texttt{Near}\-\texttt{Ring}}
\label{FactorizedGroupAssociatedWithCircleSubgroupOfNearRing}
}\hfill{\scriptsize (operation)}}\\
\textbf{\indent Returns:\ }
a group 



 The arguments are a nearring $R$ with identity and a circle subgroup $H$ of $R$. The output is the semidirect product associated with $H$ and the automorphisms determined by the corresponding units. }

 
\begin{Verbatim}[commandchars=!@|,fontsize=\small,frame=single,label=Example]
  !gapprompt@gap>| !gapinput@FG:=FactorizedGroupAssociatedWithCircleSubgroupOfNearRing(T,F[2]);|
  <pc group with 2 generators>
  !gapprompt@gap>| !gapinput@IdGroup(FG);|
  [ 25, 2 ]
\end{Verbatim}
 

\subsection{\textcolor{Chapter }{ConstantPartOfNearRing}}
\logpage{[ 2, 2, 16 ]}\nobreak
\hyperdef{L}{X790356BF873963F4}{}
{\noindent\textcolor{FuncColor}{$\triangleright$\enspace\texttt{ConstantPartOfNearRing({\mdseries\slshape R})\index{ConstantPartOfNearRing@\texttt{ConstantPartOfNearRing}}
\label{ConstantPartOfNearRing}
}\hfill{\scriptsize (attribute)}}\\
\textbf{\indent Returns:\ }
a constant part 



 The argument is a nearring $R$. The output is the constant part of the nearring $R$. }

 
\begin{Verbatim}[commandchars=!@|,fontsize=\small,frame=single,label=Example]
  !gapprompt@gap>| !gapinput@H:=LocalNearRing(361,2,342,7,7);|
  ExplicitMultiplicationNearRing ( <pc group of size 361 with 
  2 generators> , multiplication )
  !gapprompt@gap>| !gapinput@C:=ConstantPartOfNearRing(H);;|
  !gapprompt@gap>| !gapinput@Size(C);|
  19
\end{Verbatim}
 

\subsection{\textcolor{Chapter }{ZeroSymmetricPartOfNearRing}}
\logpage{[ 2, 2, 17 ]}\nobreak
\hyperdef{L}{X7F2C663584E108AD}{}
{\noindent\textcolor{FuncColor}{$\triangleright$\enspace\texttt{ZeroSymmetricPartOfNearRing({\mdseries\slshape R})\index{ZeroSymmetricPartOfNearRing@\texttt{ZeroSymmetricPartOfNearRing}}
\label{ZeroSymmetricPartOfNearRing}
}\hfill{\scriptsize (attribute)}}\\
\textbf{\indent Returns:\ }
a zero\texttt{\symbol{45}}symmetric part 



 The argument is a nearring $R$. The output is the zero\texttt{\symbol{45}}symmetric part of the nearring $R$. }

 
\begin{Verbatim}[commandchars=!@|,fontsize=\small,frame=single,label=Example]
  !gapprompt@gap>| !gapinput@ZeroSymmetricPartOfNearRing(H);;|
  !gapprompt@gap>| !gapinput@Size(last);|
  19
\end{Verbatim}
 

\subsection{\textcolor{Chapter }{GroupOfUnitsAsGroupOfAutomorphisms}}
\logpage{[ 2, 2, 18 ]}\nobreak
\hyperdef{L}{X8000A8D98567B2F8}{}
{\noindent\textcolor{FuncColor}{$\triangleright$\enspace\texttt{GroupOfUnitsAsGroupOfAutomorphisms({\mdseries\slshape R})\index{GroupOfUnitsAsGroupOfAutomorphisms@\texttt{GroupOfUnitsAsGroupOfAutomorphisms}}
\label{GroupOfUnitsAsGroupOfAutomorphisms}
}\hfill{\scriptsize (attribute)}}\\
\textbf{\indent Returns:\ }
a group 



 The argument is a nearring $R$. The output is a group of automorphisms of the additive group of $R$ that is isomorphic to the group of units of $R$. }

 
\begin{Verbatim}[commandchars=!@|,fontsize=\small,frame=single,label=Example]
  !gapprompt@gap>| !gapinput@M:=LocalNearRing(27,4,18,3,2);  |
  ExplicitMultiplicationNearRing ( <pc group of size 27 with 3 generators> , multiplication )
  !gapprompt@gap>| !gapinput@GroupOfUnitsAsGroupOfAutomorphisms(M);|
  <group of size 18 with 2 generators>
  !gapprompt@gap>| !gapinput@Size(last);|
  18                
\end{Verbatim}
 

\subsection{\textcolor{Chapter }{IsDistributiveElementOfNearRing}}
\logpage{[ 2, 2, 19 ]}\nobreak
\hyperdef{L}{X870157C878E9606B}{}
{\noindent\textcolor{FuncColor}{$\triangleright$\enspace\texttt{IsDistributiveElementOfNearRing({\mdseries\slshape R, r})\index{IsDistributiveElementOfNearRing@\texttt{IsDistributiveElementOfNearRing}}
\label{IsDistributiveElementOfNearRing}
}\hfill{\scriptsize (operation)}}\\
\textbf{\indent Returns:\ }
a boolean 



 The argument is a nearring $R$ and an element $r$. The output is \texttt{true} if $r$ is a distributive element of the nearring $R$, otherwise the output is \texttt{false}. }

 
\begin{Verbatim}[commandchars=!@|,fontsize=\small,frame=single,label=Example]
  !gapprompt@gap>| !gapinput@D:=LocalNearRing(49,2,42,6,1); |
  ExplicitMultiplicationNearRing ( <pc group of size 49 with 
  2 generators> , multiplication )
  !gapprompt@gap>| !gapinput@h:=List(D);;|
  !gapprompt@gap>| !gapinput@d:=h[3];|
  (f2^2)
  !gapprompt@gap>| !gapinput@IsDistributiveElementOfNearRing(D,d);|
  true
\end{Verbatim}
 

\subsection{\textcolor{Chapter }{IsSemiDistributiveNearRing}}
\logpage{[ 2, 2, 20 ]}\nobreak
\hyperdef{L}{X7F84066685C00F01}{}
{\noindent\textcolor{FuncColor}{$\triangleright$\enspace\texttt{IsSemiDistributiveNearRing({\mdseries\slshape R})\index{IsSemiDistributiveNearRing@\texttt{IsSemiDistributiveNearRing}}
\label{IsSemiDistributiveNearRing}
}\hfill{\scriptsize (property)}}\\
\textbf{\indent Returns:\ }
a boolean 



 The argument is a nearring $R$. The output is \texttt{true} if $R$ is a semidistributive nearring, otherwise the output is \texttt{false}. }

 
\begin{Verbatim}[commandchars=!@|,fontsize=\small,frame=single,label=Example]
  !gapprompt@gap>| !gapinput@N:=LocalNearRing(16,10,8,2,7); |
  ExplicitMultiplicationNearRing ( <pc group of size 16 with 4 generators> , multiplication )
  !gapprompt@gap>| !gapinput@IsSemiDistributiveNearRing(N);|
  true
\end{Verbatim}
 

\subsection{\textcolor{Chapter }{IsNearRingWithIdentity}}
\logpage{[ 2, 2, 21 ]}\nobreak
\hyperdef{L}{X798571F68791D67A}{}
{\noindent\textcolor{FuncColor}{$\triangleright$\enspace\texttt{IsNearRingWithIdentity({\mdseries\slshape R})\index{IsNearRingWithIdentity@\texttt{IsNearRingWithIdentity}}
\label{IsNearRingWithIdentity}
}\hfill{\scriptsize (property)}}\\
\textbf{\indent Returns:\ }
a boolean 



 The argument is a nearring $R$. The output is \texttt{true} if $R$ is a nearring with identity, otherwise the output is \texttt{false}. }

 
\begin{Verbatim}[commandchars=!@|,fontsize=\small,frame=single,label=Example]
  !gapprompt@gap>| !gapinput@N:=LocalNearRing(343,5,294,8,2);    |
  ExplicitMultiplicationNearRing ( <pc group of size 343 with 
  3 generators> , multiplication )
  !gapprompt@gap>| !gapinput@IsNearRingWithOne(N);|
  false
  !gapprompt@gap>| !gapinput@Identity(N);|
  (f1)
  !gapprompt@gap>| !gapinput@IsNearRingWithIdentity(N);|
  true
\end{Verbatim}
 

\subsection{\textcolor{Chapter }{IsSubNearRing}}
\logpage{[ 2, 2, 22 ]}\nobreak
\hyperdef{L}{X7A513FD081323B01}{}
{\noindent\textcolor{FuncColor}{$\triangleright$\enspace\texttt{IsSubNearRing({\mdseries\slshape R, H})\index{IsSubNearRing@\texttt{IsSubNearRing}}
\label{IsSubNearRing}
}\hfill{\scriptsize (operation)}}\\
\textbf{\indent Returns:\ }
a boolean 



 The arguments are a nearring $R$ with identity and a subgroup $H$ of the additive group of $R$. The output is \texttt{true} if $H$ is the additive group of a subnearring of $R$, otherwise the output is \texttt{false}. }

 
\begin{Verbatim}[commandchars=!@|,fontsize=\small,frame=single,label=Example]
  !gapprompt@gap>| !gapinput@T:=LocalNearRing(49,2,42,1,2);        |
  ExplicitMultiplicationNearRing ( <pc group of size 49 with 
  2 generators> , multiplication )
  !gapprompt@gap>| !gapinput@G:=GroupReduct(T);                    |
  <pc group of size 49 with 2 generators>
  !gapprompt@gap>| !gapinput@S:=Subgroups(G);;|
  !gapprompt@gap>| !gapinput@Size(S);|
  10
  !gapprompt@gap>| !gapinput@IsSubNearRing(T,S[3]);|
  true
  !gapprompt@gap>| !gapinput@IsSubNearRing(T,S[9]); |
  false
  !gapprompt@gap>| !gapinput@D:=SmallGroup(7,1);;|
  !gapprompt@gap>| !gapinput@IsSubNearRing(T,D);   |
  false
\end{Verbatim}
 }

 }

 \def\bibname{References\logpage{[ "Bib", 0, 0 ]}
\hyperdef{L}{X7A6F98FD85F02BFE}{}
}

\bibliographystyle{alpha}
\bibliography{LocalNR.bib}

\addcontentsline{toc}{chapter}{References}

\def\indexname{Index\logpage{[ "Ind", 0, 0 ]}
\hyperdef{L}{X83A0356F839C696F}{}
}

\cleardoublepage
\phantomsection
\addcontentsline{toc}{chapter}{Index}


\printindex

\immediate\write\pagenrlog{["Ind", 0, 0], \arabic{page},}
\newpage
\immediate\write\pagenrlog{["End"], \arabic{page}];}
\immediate\closeout\pagenrlog
\end{document}
